\section{Strategy circuits, initialization rows, and padding}\label{sec:circuits}
The halting-game finite model \eqref{eq:finite-import} is explicit but
has long computation-history labels. A small system register does not
imply a polynomial measurement circuit. The following resource bound
is sufficient for the conditional composition.

\begin{lemma}[Sequential orthogonality context]\label{lem:sequential}
For projections $E,F$ on $\C^d$, Alice can answer the context forbidding
the simultaneous $E=F=1$ outcome with the POVM
\[
 E,\qquad (I-E)F(I-E),\qquad (I-E)(I-F)(I-E).
\]
On a maximally entangled state, Bob using $E^{\mathsf T},F^{\mathsf T}$
has zero $E$ disagreement and $F$ disagreement
\begin{equation}\label{eq:sequential}
 \frac{3\tr(EF)-2\tr(EFEF)}d\leq3\norm{EF}_f^2,
 \qquad \norm A_f^2=d^{-1}\tr(A^*A).
\end{equation}
\end{lemma}
\begin{proof}
The three positive effects sum to $I$. Measure $E$ first and, after
a negative answer, measure $F$; the listed effects are exactly the
resulting measurement. Let $M_F=(I-E)F(I-E)$. The maximally entangled
trace identity with Bob's transpose gives disagreement
$\tr(M_F+F-2M_FF)/d$. Expanding,
$\tr M_F=\tr F-\tr EF$ and
$\tr(M_FF)=\tr F-2\tr EF+\tr EFEF$, proving the identity.
The last trace is nonnegative because it is the trace of $(FEF)^2$
on $\operatorname{ran}F$. Also $\norm{EF}_f^2=\tr(EF)/d$.
\end{proof}
This coherent sequential POVM replaces the compiler's projection onto
$(I-E)\operatorname{ran}F$. It changes only a constant in its loss
bound, by at most a factor $3/2$. It is not an operator-norm
approximation to that different projective measurement.

\begin{proposition}[Halting-game circuit resources]\label{prop:circuits}
For every fixed $\delta>0$ and even $N$, the finite models of
\cite[Theorem 7.1]{unprice}, with the canonical compiler blocks of its
Theorem 3.1 and Lemma~\ref{lem:sequential}, give near-winning
measurement circuits on system dimension $c_0d_N$, of loss at most
$\eps_N=C_1 16^{-N}$. Their size is at most
$2^{(1+\delta)N}\operatorname{poly}(N)$ and their work space is
$O(N^2)$ qubits. They have uniform classical descriptions generated
within the same exponential scale, allowing an arbitrarily small
fixed slack in $\delta$ for synthesis overhead. Finite-gate
implementation error in the coherent measurement is at most
$2^{-6N}$.
\end{proposition}
\begin{proof}
We describe the algorithm against the explicit source model; we use,
rather than re-prove, the compiler's algebraic soundness and finite
model identities. Its coordinate registers are $N$-bit address, fuel
and time, and $k_N=O(N)$ optional slots. The flag is
\[
 f_N(x)=\sum_{e\in K_N}\one_{\{x\equiv4^e\pmod{4^{e+1}}\}}.
\]
The residue classes are disjoint. Computing the flag, digits, masks,
fuel reset predicates, and partial counts is polynomial in $N$ once
$K_N$ is known. The count $h_N(x)=\sum_{r<x}f_N(r)$ is a sum of
residue counts, using division with remainder for each $e$; no scan
over $r<x$ is needed. Compute the maximum $N$-window count defining
$k_N$ by a $2^N\operatorname{poly}(N)$ scan. Bounded simulation
computes $K_N$ in that time using polynomial space. These are
classical preprocessing steps.

The diagonal geometric series for the source contractions $B,C$
are truncated after $O(N)$ terms. Every translated slot is determined
by the same residue counts and cyclic shifts. Address-clock dressing
is the controlled slot rotation $D=\bigoplus_x R_0^{h_N(x)}$.
Time increment, bit permutations and slot rotation have explicit
involutive factorizations. For an increment coordinate, the maps
$i\mapsto-i$ and $i\mapsto1-i$ multiply to increment; these maps are
polynomial-time reversible arithmetic.

The only less immediate permutation is the dressed $V$. In the $D^*$
gauge it increments address $x$, applying a fixed slot permutation
$P_0=R_0^{h_N(2^N)}$ only on wraparound. Given a slot string $z$,
enumerate its orbit under $P_0$, of length $s\leq k_N$; choose the
least representative $z_0$ and express $z=P_0^jz_0$. The coordinate
$i=x+2^Nj$ is a cycle coordinate modulo $2^Ns$. Apply the two
involutions above in this coordinate, convert back, and conjugate by
$D$. Orbit search is over at most $k_N$ rotations, not over an
exponential cycle. An efficiently computable involution is implemented
by compute, swap, and uncompute.

For a diagonal contraction $r$, the compiler uses
\[
 u(r)=2r-1+2i\sqrt{r(1-r)}.
\]
Realification is a single controlled rotation. Its two reflection
factors are the fixed diagonal reflection and the rotation-plane
reflection. Diagonal labels are computed reversibly, the appropriate
rotation or phase applied, and the label uncomputed. Projections use
a Boolean phase; source flips use controlled bit flips.

The expensive labels have the form
$\operatorname{Step}^{,t\bmod2^j}(\operatorname{Init}(x))$.
There are fewer than $2^N$ steps. The fixed counter machine has
constantly many counters, each at most $N/2+2^N$, hence configurations
use $O(N)$ bits. Direct computation costs
$2^N\operatorname{poly}(N)$ time and $O(N)$ space. Reversible
pebbling \cite{bennett} implements compute, copy, and uncompute with
time $2^{(1+\delta)N}\operatorname{poly}(N)$ and $O(N^2)$ work
bits, for each fixed $\delta>0$.

Precision must include square-root sensitivity at endpoints. The
uniform inequality $|u(r)-u(s)|=O(\sqrt{|r-s|})$ on $[0,1]$ follows
from the square-root inequality and the Lipschitz bound for $r(1-r)$.
Use at least $16N+O(\log N)$ fractional bits for contraction values
and discard a value $2^{-m}$ only beyond that precision. This yields
source errors $O(2^{-8N})$ after square roots. For example discarding
$r=2^{-(8N+1)}$ would instead cause error $2\sqrt r$, too large for
$2^{-6N}$. Compute angles to a comparable linear precision. A
polynomial number of register-controlled rotations is synthesized
with polynomial precision overhead \cite{dawson}; allocating
$O(N+\log N)$ synthesis bits to each rotation gives total error at
most $2^{-6N}$. Arithmetic and reversible classical gates may be
implemented exactly in a fixed universal algebraic gate set.

The remaining compiler operations have presentation-dependent constant
size: realification, selector Pauli registers, controlled source words,
and doubled reflection blocks. In the source's first conjugacy block
these have forms $\operatorname{diag}(X_i,I)$ and
$\left(\begin{smallmatrix}0&X_i\\X_i&0\end{smallmatrix}\right)$;
the second uses fixed-length products such as $ab$, $bab$, and $bc$.
They require only controlled source reflections and a constant number
of ancillary selector bits. Measure each commuting context with a
constant number of controlled-reflection tests, and use
Lemma~\ref{lem:sequential} for each orthogonality context. Forbidden
answers caused by gate approximation are mapped to a fixed allowed
assignment, changing the losing probability by at most the
implementation error. Every context has an allowed assignment by the
imported compiler. The exact source defects $O(4^{-N})$ give loss
$O(16^{-N})$; adding $O(2^{-6N})$ leaves that scale unchanged.
All preprocessing and gate-label generation are covered by the stated
bound, with slack absorbing fixed polynomial factors.
\end{proof}

This is an algorithmic resource construction. It is not a supplied
executable implementation of the complete compiler, and does not
establish polynomial-size measurement circuits. No classical circuit
lower bound is inferred from the long history labels.

\subsection{Distributed initialization: a conditional algebraic estimate}
To specify what a physical layout must preserve, consider two clocks
and a row $R_j=\{t_A=j\}$. Let $u_{R_j}$ be uniform on that row and
$u$ uniform on the grid. Spectral separation in the Bob coordinate gives
\begin{equation}\label{eq:rowresistance}
 (u_{R_j}-u)^*L_2^+(u_{R_j}-u)
       =(L_{
       \rm path}^+)_{jj}/n\leq1/3.
\end{equation}
For completeness the centered path diagonal is
\[
 (L_{\rm path}^+)_{jj}
 =\frac1{n^2}\left[\sum_{i=1}^{j}i^2+\sum_{i=1}^{n-1-j}i^2\right],
\]
the squared edge-flow sum for injection at $j$ and uniform withdrawal.
In particular the corner diagonal maximizes it and equals
$(n-1)(2n-1)/(6n)<n/3$.

Suppose a proposed layout supplies $q$ commuting rewound ancilla
projections $P_j$, with a unit penalty on every vertex of their
initialization row. Jensen's inequality and positivity give
$\norm{(I-P_j)\operatorname{avg}_{R_j}\psi}\leq\sqrt{\lambda/n}$.
Together with \eqref{eq:rowresistance},
\[
 \norm{(I-P_j)m}\leq(1/\sqrt n+1/\sqrt3)\sqrt\lambda.
\]
Since $I-\prod_jP_j\leq\sum_j(I-P_j)$, initialization error for all
ancillas is at most $q(1/\sqrt n+1/\sqrt3)^2\lambda$. Include a
commuting corner question projection and set
\[
 B=(1+\sqrt{G_2})^2+q(1/\sqrt n+1/\sqrt3)^2.
\]
The proof of \eqref{eq:extraction-proof} now gives the safe coefficient
\begin{equation}\label{eq:rowconstant}
 C'=2(1+2\sqrt B+2\sqrt{G_2})^2+1/2=O(q+G_2),
\end{equation}
and the normalized reference distance is at most
$2(1/2+\sqrt B)\sqrt{M\lambda}$. These estimates are conditional
on exactly the indicated rewound row penalties; they do not build them.

\begin{lemma}[Padding]\label{lem:padding}
Idle padding to $T\geq h^{3/2}$ makes the initialized variational
energy at most $h^{-3}$ for two clocks. If also $T\leq h^C$ for a
fixed $C$, then $G_2=O(\log h)$.
\end{lemma}
\begin{proof}
The variational input has loss at most one and energy at most $1/M$,
where $M=(T+1)^2>h^3$. The upper bound on $T$ and
\eqref{eq:greenbound} give the resistance estimate. A lower bound
on $T$ alone would not imply the latter conclusion.
\end{proof}
